\subsection{Schaltverhalten eines Transistors}
% Alt: Aufgabe 14
\Aufgabe{
    \label{Aufg14}
    Gegeben ist die abgebildete Transistorschaltung mit den folgenden Parametern:
    $U_0=\mathrm{5\,V}$ und $U_L=\mathrm{15\,V}$.
    Damit der gezeigte Siliziumtransistor schaltet, benötigt er einen Basisstrom von \( I_\mathrm{B}=\mathrm{50\,\mu A} \), um in den leitenden Zustand zu gelangen.

    \begin{figure}[H]
        \centering
        \input{Tikz/src/FigTransistorSchaltverhalten.tex} \alttext{Die Abbildung zeigt eine Transistorschaltung. Links befindet sich eine Spannungsquelle, die eine Spannung \(U_0\) liefert. Von dieser Quelle aus fließt der Strom durch einen Vorwiderstand, der als \(R_\mathrm{V}\) gekennzeichnet ist weiter zur Basis eines Transistors. Der Basisstrom ist in rot als Pfeil in Stromrichtung mit \(I_\mathrm{B}\) gekennzeichnet.
            Auf der rechten Seite des Transistors ist ein Lastwiderstand \(R_\mathrm{C}\) angeschlossen, hinter dem sich eine weitere Spannungsquelle befindet die mit \(U_\mathrm{L}\) gekennzeichnet ist.}
        \label{fig:FigTransistorSchaltverhalten}
    \end{figure}

    \begin{itemize}
        \item[a)] Wie groß muss der Vorwiderstand \( R_\mathrm{V} \) gewählt werden, damit der Transistor durchschaltet?
        \item[b)] Der Transistor hat einen Stromverstärkungsfaktor \( \beta = 180 \) und einen Kollektorwiderstand \( R_\mathrm{C} = \mathrm{200\,\Omega} \).
              Wie groß sind der Kollektorstrom \( I_\mathrm{C} \) und der Spannungsabfall zwischen Kollektor und Emitter \( U_\mathrm{CE} \)?
    \end{itemize}

}


\Loesung{
    \begin{itemize}
        \item[a)]
              Berechnung des Vorwiderstands für den Basisstrom:
              \begin{equation*}
                  R_\mathrm{V}  = \frac{U}{I} = \frac{\mathrm{5\,V} - \mathrm{0{,}6\,V}}{\mathrm{50\,\mu A}} = \frac{\mathrm{4{,}4\,V}}{\mathrm{50\,\mu A}} = \mathrm{88\,k\Omega}
              \end{equation*}

        \item[b)]
              Berechnung des Kollektorstroms mit Stromverstärkungsfaktor:
              \begin{gather*}
                  \beta = \frac{I_\mathrm{C}}{I_\mathrm{B}} \\
                  I_\mathrm{C} = \beta \cdot I_\mathrm{B} = 180 \cdot \mathrm{50\,\mu A} = \mathrm{9\,mA}
              \end{gather*}

              Berechnung des Spannungsabfalls am Kollektorwiderstand:
              \begin{equation*}
                  U_\mathrm{RC} = I_\mathrm{C} \cdot R_\mathrm{C} = \mathrm{9\,mA} \cdot \mathrm{200\,\Omega} = \mathrm{1{,}8\,V}
              \end{equation*}

              Berechnung der Spannung zwischen Kollektor und Emitter:
              \begin{equation*}
                  U_\mathrm{CE} = \mathrm{15\,V} - \mathrm{1{,}8\,V} = \mathrm{13{,}2\,V}
              \end{equation*}
    \end{itemize}

    Siehe: Abschnitt \ref{sec:Bipolartransistor} und \ref{fig:StroemeUndSpannungenBeimBipolartransistor}

}